\subsection{Operações em curvas e homotopias}

Para fins de conveniência, utilizamos a seguinte notação para curvas fechadas:%
\linebreak\vspace{-1.2em}

\vspace{0.5em}

\begin{convention}
    Sejam $p$ e $q$ pontos quaisquer em $X$. Denotamos simbolicamente por $\gamma  : p \pathto q$ o juízo de que $\gamma$ é uma curva em $X$ de $p$ até $q$, i.e.\@ $\gamma : I \to X$ é uma função contínua de $I = [0, 1]$ em $X$ tal que $\gamma(0) = p$ e $\gamma(1) = q$, ou apenas por $p \pathto q$ caso $\gamma$ esteja implícita ou puder ser inferida pelo contexto.
    \qedalt
\end{convention}

Pelas observações da seção anterior, se $f$ e $g$ são funções contínuas de $X$ em $Y$, então toda homotopia $\eta : f \sim g$ pode ser unicamente representada por uma curva $\gamma : f \pathto g$ em $C(X, Y)$, dada por $\gamma(t) : x \mapsto \eta(x, t)$. Por um abuso de notação, podedmos identificar as homotopias $f \sim g$ com as curvas $f \pathto g$ que lhe representam, e em particular no desenvolvimento da seção presente, os enunciados para curvas contêm as homotopias como um caso particular.

Há uma operação muito natural entre curvas, denominada \emph{concatenação}, e que intuitivamente consiste em ``colar'' duas curvas $\alpha : p \pathto q$ e $\beta : q \pathto r$ em um mesmo espaço $X$ que compartilham uma mesma extremidade $q \in X$. O resultado é uma curva $\alpha \cdot \beta : p \pathto r$ entre as extremidades restantes e que passa por $q$ no caminho.%
\linebreak\vspace{-1.2em}

\begin{figure}[htpb]
\centering
\begin{tikzpicture}
    \draw [thick, dotted]
        (-2,-1)
        to[out=45,in=135]
            node [above] {$\alpha$}
        (0,-1)
        ;
    
    \draw [thick, densely dashed]
        (0,-1)
        to[out=-45,in=-90]
            node [below] {$\beta$}
        (2,0)
        ;

    \node
        [left] at (-2,-1) {$p$};
    \node
        [below left] at (0,-1) {$q$};
    \node
        [right] at (2,0) {$r$};

    \draw [fill]
        (-2,-1) circle (1pt)
        (0,-1) circle (1pt)
        (2,0) circle (1pt)
        ;

    \node at (3,-0.5) {$\mapsto$};

    \draw [thick]
        (4,-1)
        to[out=45,in=135]
        (6,-1)
            node [above right] {$\alpha \cdot \beta$}
        to[out=-45,in=-90]
        (8,0)
        ;

    \node
        [left] at (4,-1) {$p$};
    \node
        [below left] at (6,-1) {$q$};
    \node
        [right] at (8,0) {$r$};

    \draw [fill]
        (4,-1) circle (1pt)
        (6,-1) circle (1pt)
        (8,0) circle (1pt)
        ;
\end{tikzpicture}
\end{figure}

Podemos definir $\alpha \cdot \beta$ como qualquer uma entre várias curvas que capturam esta intuição e em adição são homotópicas entre si, mas uma definição bastante natural é a seguinte. {\color{red} TODO: intuição da concatenação}

\vspace{0.5em}

\begin{definition}[(Concatenação)]
    Sejam $p$, $q$ e $r$ pontos em $X$, e sejam $\alpha : p \pathto q$ e $\beta : q \pathto r$ curvas em $X$. Denotamos por $\alpha \cdot \beta : p \pathto r$ uma curva em $X$, dada por
%
    \begin{equation*}
        (\alpha \cdot \beta)(t) =
        \begin{cases}
            \alpha(2t) & \text{se $0 \leq t \leq 0,5$} \\
            \beta(2t - 1) & \text{se $0,5 \leq t \leq 1$}
        \end{cases}
    \end{equation*}
%
    para todo $t \in I$, e denominada a \emph{concatenação} de $\alpha$ e $\beta$.
    \qedalt
\end{definition}

{\color{red} TODO: intuição da curva inversa}

\vspace{0.5em}

\begin{definition}[(Curva inversa)]
    Sejam $p$ e $q$ pontos em $X$, e seja $\gamma : p \pathto q$ uma curva em $X$. Denotamos por $\gamma^{-1} : q \pathto p$ uma curva em $X$, a \emph{inversa} de $\gamma$, dada por%
    \linebreak\vspace{-0.6em}
%
    \begin{equation*}
        \gamma^{-1}(t) = \gamma(1 - t)
    \end{equation*}
%
    para todo $t \in I$.
    \qedalt
\end{definition}

{\color{red} TODO: intuição da curva constante}

\vspace{0.5em}

\begin{definition}[(Curva constante)]
    Dado um ponto qualquer $p$ em $X$, definimos por
%
    \begin{equation*}
        \msf{const}_p(t) = p
    \end{equation*}
%
    para todo $t \in I$ a \emph{curva constante em $p$}.
    \qedalt
\end{definition}

{\color{red} TODO: conexão por curvas, componentes conexas por curvas}

\vspace{0.5em}

\begin{proposition}
    $\pathto$ é uma relação de equivalência sobre $X$. Mais especificamente:
%
    \begin{enumerate}
    [
        itemsep=0pt,
        label={(\alph*)}
    ]
    
        \item 
        $\msf{const}_p : p \pathto p$
        \hfill
        ($p \in X$)

        \item 
        $\gamma : p \pathto q$ $\implies$ $\gamma^{-1} : q \pathto p$
        \hfill
        ($p, q \in X$)

        \item 
        $\alpha : p \pathto q$, $\beta : q \pathto r$ $\implies$ $\alpha \cdot \beta : p \pathto q$
        \hfill
        ($p, q, r \in X$)
    \end{enumerate}
%
    Logo, as componentes conexas por curvas em $X$ formam uma partição de $X$.
    \qedalt
\end{proposition}

{\color{red} TODO: homotopias representam curvas}

\vspace{0.5em}

\begin{corollary}
    $\sim$ é uma relação de equivalência sobre $C(X,Y)$. Em particular, se identificarmos uma homotopia $\eta : f \sim g$ entre $f, g \in C(X, Y)$ com a curva $f \pathto g$ em $C(X, Y)$ que representa $\eta$, então mais especificamente:

    \begin{enumerate}
    [
        itemsep=0pt,
        label={(\alph*)}
    ]
    
        \item 
        $\msf{const}_f : f \sim f$
        \hfill
        ($f \in C(X, Y)$)

        \item 
        $\eta : f \sim g$ $\implies$ $\eta^{-1} : g \sim f$
        \hfill
        ($f, g \in C(X, Y)$)

        \item 
        $\alpha : f \sim g$, $\beta : g \sim h$ $\implies$ $\alpha \cdot \beta : f \sim h$
        \hfill
        ($f, g, h \in C(X, Y)$)
    \end{enumerate}
%
    Logo, as classes de homotopia em $C(X,Y)$ formam uma partição de $C(X,Y)$.
    \qedalt
\end{corollary}